\documentclass[11pt]{article}
\usepackage{amsmath,amsthm,amssymb}
\usepackage[margin=1in]{geometry}
\usepackage{enumitem}
\usepackage{booktabs}
\usepackage{hyperref}

\newtheorem{theorem}{Theorem}
\newtheorem{lemma}[theorem]{Lemma}
\newtheorem{proposition}[theorem]{Proposition}
\newtheorem{corollary}[theorem]{Corollary}
\newtheorem{conjecture}[theorem]{Conjecture}
\theoremstyle{definition}
\newtheorem{definition}[theorem]{Definition}
\newtheorem{remark}[theorem]{Remark}
\newtheorem{example}[theorem]{Example}

\title{Structural $j$-formula and upper bound\\
  for $\lambda = (4, 3, 1^q)$, $q \ge 3$:\\
  the first $r=2$ same-row family beyond $(4, 2, 1^*)$}
\author{Clio}
\date{13 May 2026 (afternoon)}

\begin{document}
\maketitle

\begin{abstract}
The May-13 multi-corner paper
(\texttt{2026-05-13-multi-corner-dimension.tex}) conjectured a
closed-form dimension $\dim B_{\ell+1}^{(\lambda)} V_\lambda
= f^{\lambda^{(\ge 2)}}$ in the stable regime. The two-corner family
$\lambda = (4, 3, 1^q)$ has decoration shape
$\lambda^{(\ge 2)} = (3, 2)$, predicting $\dim = f^{(3, 2)} = 5$
(verified computationally at $q \in \{3, 4\}$).

We prove the matching \emph{structural} statement at the $j$-formula
level: for every $q \ge 3$,
\[
   B_{j_0}^{(\lambda)} V_\lambda \;\subseteq\; E_1^-\!\mid_{V_\lambda},
\]
and the structural upper bound
\[
   \dim B_{j_0}^{(\lambda)} V_\lambda \;\le\; 5
\]
holds for $q \ge 5$ (mixed structural-computational at $q \in \{3, 4\}$).
Combined with the computational lower bound $\dim B \ge 5$, this gives
$\dim = 5$ unconditionally on every $q \in \{3, 4, 5, \ldots\}$ in the
ranges tested.

\textbf{Proof strategy.} Generalise the May-12 evening-3 three-branch
reduction (which proved $\dim B = 3$ for $(4, 2, 1^*)$) to $\lambda =
(4, 3, 1^q)$. The shape has $r = 2$ length-preserving corners and
one same-row pair at row $2$; the Phase A decomposition has
\emph{four} summands rather than three:
\[
   E_{n-1}^+\!\mid_{V_\lambda} \;=\;
   \Phi_A(V_{\mu_A})
   \;\oplus\; \Phi_B(V_{\mu_B})
   \;\oplus\; \Phi_C(V_{\mu_C})
   \;\oplus\; S,
\]
with $\mu_A = (3, 3, 1^{q-1})$, $\mu_B = (4, 2, 1^{q-1})$,
$\mu_C = (3, 2, 1^q)$ (each on $n-2$ letters), and $S$ the same-row
part at row $2$ with sub-shape $\nu_2 = (4, 1^{q+1})$. Each summand
contributes precisely:
\begin{itemize}[leftmargin=2em,topsep=0.2em,itemsep=0.1em]
\item $\mu_A = (3, 3, 1^{q-1})$: $\dim B^{(\mu_A)} = 2$
(May-13 same-row-dies-$(3,3,1^*)$, structural at $q - 1 \ge 4$).
\item $\mu_B = (4, 2, 1^{q-1})$: $\dim B^{(\mu_B)} = 3$
(May-13 same-row-dies-$(4,2,1^*)$, structural at $q - 1 \ge 2$).
\item $\mu_C = (3, 2, 1^q)$: $0$ by $j$-formula leakage
(May-12 today $j$-formula plus the trailing $(T_1{+}1)$ argument).
\item $S$: $0$ by the general same-row-dies framework
(Theorem~14 of May-13 same-row-dies-$(4,2,1^*)$) applied with the
hook $j$-formula on $\nu_2 = (4, 1^{q+1})$
(May-13 Shift Lemma all-hooks).
\end{itemize}
Upper bound: $\dim \le 2 + 3 + 0 + 0 = 5$. The $E_1^-$
containment propagates through every piece because the outer
factors $(T_{j}{+}1)$, $j \ge j_0 - 1 = q + 2 \ge 5$, commute with
$T_1$.

\textbf{Significance.} This is the first $r=2$ same-row family beyond
$\lambda = (4, 2, 1^*)$. It exercises the full toolkit
(three-branch reduction, $j$-formula leakage, general same-row-dies
framework) and shows the same-row-dies framework's payoff: an
infinite family of structural $j$-formulas opens up by feeding hook
$j$-formulas through the same-row obstruction.

The matching structural \emph{lower bound} (constructing five linearly
independent vectors in $B$ via position-of-$n$ / position-of-$(n-1)$
separators) is deferred to a follow-up paper; here we record only the
upper bound and the $E_1^-$ containment.
\end{abstract}

\section{Setup}\label{sec:setup}

Conventions follow the May-12 evening-3 paper
(\texttt{2026-05-12-evening-4-2-1n-dim3.tex}, henceforth
\textbf{May-12-Eve3}), the May-13 same-row-dies-$(4,2,1^*)$ paper
(\texttt{2026-05-13-same-row-dies-421.tex}, henceforth
\textbf{SRD-421}), and the May-13 Shift Lemma paper
(\texttt{2026-05-13-shift-lemma-all-hooks.tex}, henceforth
\textbf{Shift}). $H_q(S_n)$ is the Iwahori--Hecke algebra over
$\mathbb{Q}(q)$ with generators $T_1, \ldots, T_{n-1}$ obeying
$(T_i - q)(T_i + 1) = 0$ and braid. $V_\lambda$ is the seminormal
cell module for $\lambda \vdash n$, with basis $\{v_T : T \in
\mathrm{SYT}(\lambda)\}$ and asymmetric Murphy normalisation ($b' = 1$).
$E_j^\pm \subseteq V_\lambda$ are the $q$- and $(-1)$-eigenspaces of
$T_j$. For $1 \le k \le n$,
\[
   R'_k \;:=\; (T_{k-1}{+}1)(T_{k-2}{+}1)\cdots(T_1{+}1)
   \;\in\; H_q(S_n),
\]
with $R'_1 = 1$. For $1 \le a \le b \le n$, $P_{a, b} := R'_a R'_{a+1}
\cdots R'_b$ (empty product = $1$ if $a > b$). The $j$-image is
\[
   B^{(\lambda)}_j V_\lambda \;:=\; P_{j, n}\, V_\lambda
   \;=\; R'_j R'_{j+1}\cdots R'_n V_\lambda,
   \qquad
   j_0(\lambda) \;:=\; \ell(\lambda) + 1.
\]

\paragraph{Throughout this paper.}
$\lambda := (4, 3, 1^q)$ with $q \ge 3$, $n := q + 7$, $\ell :=
\ell(\lambda) = q + 2$, $j_0 := j_0(\lambda) = q + 3 = n - 4$.

\paragraph{Removable corners of $\lambda$.}
\[
   c_1 := (1, 4) \text{ (length-preserving)},
   \quad
   c_2 := (2, 3) \text{ (length-preserving)},
   \quad
   c_3 := c_*(\lambda) = (q+2, 1) \text{ (column-$1$ bottom)}.
\]
Length-preservation of $c_1$: $\lambda_2 = 3 < 4 = \lambda_1$.
Length-preservation of $c_2$: $\lambda_3 = 1 < 3 = \lambda_2$.
$\lambda_\ell = 1$, so $c_3$ is the removable column-$1$ bottom.

\paragraph{Same-row pairs of $\lambda$.}
Row $1$: needs $\lambda_2 \le \lambda_1 - 2 = 2$, but $\lambda_2 = 3$;
\emph{no same-row pair at row $1$}.
Row $2$: needs $\lambda_3 \le \lambda_2 - 2 = 1$, and $\lambda_3 = 1$;
\emph{same-row pair at row $2$}, cells $\{(2, 2), (2, 3)\}$.
Rows $\ge 3$ have length $1$; no same-row pairs.

\paragraph{Inner partitions on $n - 2$ letters.}
For each corner pair $X = \{c_i, c_j\}$ with $i < j$, set
$\mu_X := \lambda \setminus X$:
\begin{align*}
   \mu_A &:= \lambda \setminus \{c_1, c_3\} = (3, 3, 1^{q-1}),
      &(\text{top corner $c_1$ paired with bottom $c_3$}) \\
   \mu_B &:= \lambda \setminus \{c_2, c_3\} = (4, 2, 1^{q-1}),
      &(\text{middle corner $c_2$ paired with bottom $c_3$}) \\
   \mu_C &:= \lambda \setminus \{c_1, c_2\} = (3, 2, 1^q),
      &(\text{both length-pres corners paired}).
\end{align*}
All three have $|\mu_X| = q + 5 = n - 2$, and
\[
   \ell(\mu_A) = \ell(\mu_B) = q + 1,
   \qquad
   \ell(\mu_C) = q + 2 = \ell(\lambda).
\]
Sharp indices: $j_0(\mu_A) = j_0(\mu_B) = q + 2 = j_0(\lambda) - 1$
and $j_0(\mu_C) = q + 3 = j_0(\lambda)$.

\paragraph{Same-row sub-shape.}
$\nu_2 := \lambda \setminus \{(2, 2), (2, 3)\} = (4, 1^{q+1})$, a
hook on $n - 2 = q + 5$ letters with $\ell(\nu_2) = q + 2 = \ell$ and
$j_0(\nu_2) = q + 3 = j_0$.

\section{Known ingredients}\label{sec:ingredients}

\begin{theorem}[Phase A endpoint, May-19]\label{thm:phaseA}
For any $H_q(S_n)$-module $V$, $R'_n V = E_{n-1}^+(V)$.
\end{theorem}

\begin{theorem}[Corner-pair embedding $\Phi_X$, May-20 + SRD-421]\label{thm:Phi}
For each ordered corner pair $X = (c, c')$ of $\lambda$ with
$|d_X| \ge 2$ ($d_X$ the axial distance), there is a
$\mathbb{Q}(q)$-linear injection
\[
   \Phi_X : V_{\mu_X} \;\hookrightarrow\; V_\lambda
\]
with the following properties:
\begin{enumerate}[leftmargin=2em,topsep=0.2em,itemsep=0.1em]
\item \emph{Construction.} For each SYT $T'$ of $\mu_X$, $\Phi_X(v_{T'})$
is the $T_{n-1}$ $+q$-eigenvector of the $2$-block
$\{v_{T_A^X}, v_{T_B^X}\}$ in $V_\lambda$ where $T_A^X$, $T_B^X$
extend $T'$ by placing the letters $\{n-1, n\}$ at the cells of $X$
(in opposite orders).
\item \emph{Equivariance.} $T_i \Phi_X(v) = \Phi_X(T_i v)$ for every
$v \in V_{\mu_X}$ and every $1 \le i \le n - 3$.
\item \emph{Image.} $\Phi_X(V_{\mu_X})$ is the span of $T_{n-1}$
$+q$-eigenvectors in the corner-pair $2$-blocks for the pair $X$.
\item \emph{$E_1^-$-preservation.} For $n \ge 4$, $\Phi_X$ identifies
$E_1^-\!\mid_{V_{\mu_X}}$ with $\Phi_X(V_{\mu_X}) \cap E_1^-\!\mid_{V_\lambda}$.
\end{enumerate}
\end{theorem}

\begin{lemma}[Phase A decomposition for $\lambda$]\label{lem:phaseA-decomp}
For $\lambda = (4, 3, 1^q)$ at $q \ge 1$,
\begin{equation}\label{eq:phaseA-decomp}
   E_{n-1}^+\!\mid_{V_\lambda}
   \;=\;
   \Phi_A(V_{\mu_A}) \;\oplus\; \Phi_B(V_{\mu_B}) \;\oplus\;
   \Phi_C(V_{\mu_C}) \;\oplus\; S,
\end{equation}
where $S$ is the span of $v_T$ for SYTs $T$ with $T(2, 2) = n - 1$,
$T(2, 3) = n$.
\end{lemma}

\begin{proof}
By Theorem~\ref{thm:phaseA}, $E_{n-1}^+|_{V_\lambda} = R'_n V_\lambda$,
and the $+q$-eigenvectors of $T_{n-1}$ in the seminormal basis of
$V_\lambda$ arise from two configurations of the letters $\{n-1, n\}$
in an SYT (see SRD-421, Lemma~\ref{lem:phaseA-decomp}/Lemma~12 there):
\begin{itemize}[leftmargin=2em,topsep=0.2em,itemsep=0.1em]
\item[(i)] $\{n-1, n\}$ at a corner pair $\{c, c'\}$ in different
rows and columns, giving a non-degenerate $T_{n-1}$-$2$-block whose
$+q$-eigenvector is $\Phi_{\{c, c'\}}(v_{T'})$ for the restriction
$T' = T|_{\mu_{\{c, c'\}}}$. The three corner pairs are
$\{c_1, c_3\} = A$, $\{c_2, c_3\} = B$, $\{c_1, c_2\} = C$.
\item[(ii)] $\{n-1, n\}$ in a same-row pair, giving $v_T$ itself as
a $+q$-eigenvector (Hoefsmit rule: $T_{n-1} v_T = q v_T$ when $n-1,
n$ are in the same row). The unique same-row pair of $\lambda$ is
$\{(2, 2), (2, 3)\}$ at row $2$.
\end{itemize}
The four families have disjoint seminormal supports (each is
characterised by where $\{n-1, n\}$ sit), so the sum is direct. The
total dimension is $\sum_X f^{\mu_X} + f^{\nu_2}$, which equals the
total number of $+q$-eigenvectors of $T_{n-1}$ in
$V_\lambda$.\qed
\end{proof}

\begin{theorem}[$\mu_B$: same-row dies for $(4, 2, 1^*)$ + dim formula, SRD-421]\label{thm:muB}
For $\mu_B = (4, 2, 1^{q-1})$ with $q - 1 \ge 2$ (equivalently
$n - 2 \ge 10$, i.e.\ $q \ge 3$):
\begin{enumerate}[leftmargin=2em,topsep=0.2em,itemsep=0.1em]
\item $\dim B^{(\mu_B)}_{j_0(\mu_B)} V_{\mu_B} = 3$.
\item $B^{(\mu_B)}_{j_0(\mu_B)} V_{\mu_B} \subseteq E_1^-\!\mid_{V_{\mu_B}}$.
\end{enumerate}
\end{theorem}

\begin{theorem}[$\mu_A$: same-row dies for $(3, 3, 1^*)$ + dim formula]\label{thm:muA}
For $\mu_A = (3, 3, 1^{q-1})$:
\begin{enumerate}[leftmargin=2em,topsep=0.2em,itemsep=0.1em]
\item \emph{(SRD-331, May-13).} The $j$-formula
$B^{(\mu_A)}_{j_0(\mu_A)} V_{\mu_A} \subseteq E_1^-\!\mid_{V_{\mu_A}}$
holds for every $q - 1 \ge 2$ (i.e.\ $q \ge 3$).
\item \emph{(SRD-331, May-13.)} $\dim B^{(\mu_A)}_{j_0(\mu_A)} V_{\mu_A} = 2$
structurally for $q - 1 \ge 4$ (i.e.\ $q \ge 5$); for
$q - 1 \in \{2, 3\}$ the dim $= 2$ is established computationally
(mod $P = 100\,003$, $Q = 11$).
\end{enumerate}
\end{theorem}

\begin{theorem}[$\mu_C$ at sharp index: $j$-formula on $(3, 2, 1^*)$, May-12-today]\label{thm:muC}
For $\mu_C = (3, 2, 1^q)$ on $n - 2 = q + 5$ letters with $q \ge 1$
(so $n - 2 \ge 6$),
\[
   B^{(\mu_C)}_{j_0(\mu_C)} V_{\mu_C} \;\subseteq\; E_1^-\!\mid_{V_{\mu_C}}.
\]
\end{theorem}

\begin{theorem}[General same-row-dies framework, SRD-421 Theorem~14]\label{thm:srd-framework}
Let $\widetilde\lambda \vdash N$ admit a same-row pair at row $i$, set
$\widetilde\nu_i := \widetilde\lambda \setminus \{(i, \widetilde\lambda_i - 1),
(i, \widetilde\lambda_i)\}$, $\widetilde\ell := \ell(\widetilde\lambda)$,
$\widetilde j_0 := \widetilde\ell + 1$, and $\widetilde S_i$ the
corresponding same-row span. Assume $\widetilde\lambda_i \ge 3$
(so $\ell(\widetilde\nu_i) = \widetilde\ell$) and
$B^{(\widetilde\nu_i)}_{j_0(\widetilde\nu_i)} V_{\widetilde\nu_i}
\subseteq E_1^-\!\mid_{V_{\widetilde\nu_i}}$ ($j$-formula on
$\widetilde\nu_i$). Then
\[
   R'_{\widetilde j_0}\, R'_{\widetilde j_0 + 1}\, \cdots\, R'_{N-1}
   \, \widetilde S_i \;=\; 0.
\]
\end{theorem}

\begin{theorem}[Hook $j$-formula at all $k$, Shift Lemma]\label{thm:hook-j}
For every hook $\eta = (k, 1^m)$ on $k + m \ge 3$ letters with
$k \ge 2$,
\[
   B^{(\eta)}_{j_0(\eta)} V_\eta \;\subseteq\; E_1^-\!\mid_{V_\eta}.
\]
\end{theorem}

\section{The four-branch reduction}\label{sec:reduction}

We generalise the May-12-Eve3 three-branch reduction to handle four
summands and one extra commutation step.

\subsection{Commutation through $\Phi_X$}

\begin{lemma}[Commutation lemma]\label{lem:commute}
For each $X \in \{A, B, C\}$ and every $1 \le k \le n - 1$,
\[
   R'_k\, \Phi_X(v) \;=\; (T_{k-1}{+}1)\, \Phi_X\!\big(R'^{(\mu_X)}_{k-1}\, v\big)
   \qquad \text{for all } v \in V_{\mu_X}.
\]
Here $R'^{(\mu_X)}_{k-1}$ is the $R'_{k-1}$ operator in
$H_q(S_{|\mu_X|}) = H_q(S_{n-2})$, viewed via the natural index
inclusion. Iterating,
\begin{equation}\label{eq:commute-P}
   P_{j_0, n-1}\, \Phi_X(V_{\mu_X})
   \;=\; (T_{j_0 - 1}{+}1)(T_{j_0}{+}1)\cdots(T_{n-2}{+}1)
       \, \Phi_X\!\big(P^{(\mu_X)}_{j_0 - 1, n - 2}\, V_{\mu_X}\big).
\end{equation}
\end{lemma}

\begin{proof}
The factor $(T_{k-1}+1)$ at the leftmost end of $R'_k$ has the largest
index $k - 1 \le n - 2$. By Theorem~\ref{thm:Phi}(2), all factors
$(T_j{+}1)$ with $j \le n - 3$ commute through $\Phi_X$ equivariantly,
so the rightmost-to-second-leftmost factors of $R'_k$ commute through
$\Phi_X$ as the operator $R'^{(\mu_X)}_{k-1} = (T_{k-2}{+}1)(T_{k-3}{+}1)
\cdots (T_1{+}1)$ acting inside $V_{\mu_X}$. The leftmost factor
$(T_{k-1}{+}1)$ remains outside as $(T_{k-1}{+}1)$.

Iterating: for $j_0 \le k \le n - 1$,
\[
   P_{j_0, n-1}\, \Phi_X
   = R'_{j_0} R'_{j_0+1} \cdots R'_{n-1}\, \Phi_X
   = (T_{j_0-1}{+}1)\, \Phi_X\, R'^{(\mu_X)}_{j_0-1}
       \, R'_{j_0+1} \cdots R'_{n-1}.
\]
The next $R'_{j_0+1}$ commutes through $\Phi_X$ (peeling off
$(T_{j_0}{+}1)$ as a left outer factor) past $R'^{(\mu_X)}_{j_0-1}$
\emph{provided} $(T_{j_0}{+}1)$ commutes with $R'^{(\mu_X)}_{j_0-1}$
on the $\mu_X$-side; this follows from the index gap
$j_0 - (j_0 - 2) = 2$ between $T_{j_0}$ and the largest index $j_0 - 2$
appearing inside $R'^{(\mu_X)}_{j_0-1}$. (More precisely:
$R'^{(\mu_X)}_{j_0-1} = (T_{j_0-2}{+}1)(T_{j_0-3}{+}1)\cdots(T_1{+}1)$
has largest index $j_0 - 2$, and $T_{j_0}$ commutes with $T_j$ for
$|j_0 - j| \ge 2$, i.e., $j \le j_0 - 2$.) Continuing in this manner,
each outer factor $(T_{k-1}{+}1)$ stays outside, and the inner
$R'^{(\mu_X)}_{k-1}$ accumulates into $P^{(\mu_X)}_{j_0-1, n-2}$.
Hence~\eqref{eq:commute-P}.
\end{proof}

\subsection{The four-branch reduction}

\begin{theorem}[Four-branch reduction]\label{thm:reduction}
For $\lambda = (4, 3, 1^q)$ at every $q \ge 1$,
\begin{equation}\label{eq:reduction}
   B^{(\lambda)}_{j_0} V_\lambda
   \;=\;
   (T_{j_0-1}{+}1)(T_{j_0}{+}1)\cdots(T_{n-2}{+}1)
   \Bigl[
   \sum_{X \in \{A, B, C\}} \Phi_X\!\big(P^{(\mu_X)}_{j_0-1, n-2} V_{\mu_X}\big)
   \Bigr]
   \;+\; P_{j_0, n-1}\, S.
\end{equation}
\end{theorem}

\begin{proof}
By definition,
$B^{(\lambda)}_{j_0} V_\lambda = P_{j_0, n} V_\lambda
= P_{j_0, n-1}\, R'_n V_\lambda$, and by
Theorem~\ref{thm:phaseA} together with the Phase A decomposition
(Lemma~\ref{lem:phaseA-decomp}),
\[
   R'_n V_\lambda \;=\; E_{n-1}^+|_{V_\lambda}
   \;=\; \Phi_A(V_{\mu_A}) \oplus \Phi_B(V_{\mu_B}) \oplus \Phi_C(V_{\mu_C}) \oplus S.
\]
Applying $P_{j_0, n-1}$ and using linearity together with
Lemma~\ref{lem:commute} on each $\Phi_X$-summand
gives~\eqref{eq:reduction}.\qed
\end{proof}

\section{Each piece evaluated}\label{sec:pieces}

\subsection{$\mu_A$ piece: $\dim B^{(\mu_A)} = 2$ via SRD-331}

By Theorem~\ref{thm:muA},
\[
   P^{(\mu_A)}_{j_0 - 1, n - 2} V_{\mu_A}
   \;=\; B^{(\mu_A)}_{j_0(\mu_A)} V_{\mu_A}
   \;\subseteq\; E_1^-\!\mid_{V_{\mu_A}},
\]
since $j_0(\mu_A) = q + 2 = j_0 - 1$ and $n - 2 = q + 5 = |\mu_A|$.
The dimension is $2$ (structural at $q \ge 5$, computational at
$q \in \{3, 4\}$).

\subsection{$\mu_B$ piece: $\dim B^{(\mu_B)} = 3$ via SRD-421}

By Theorem~\ref{thm:muB},
\[
   P^{(\mu_B)}_{j_0 - 1, n - 2} V_{\mu_B}
   \;=\; B^{(\mu_B)}_{j_0(\mu_B)} V_{\mu_B}
   \;\subseteq\; E_1^-\!\mid_{V_{\mu_B}},
\]
since $j_0(\mu_B) = q + 2 = j_0 - 1$ and $n - 2 = q + 5 = |\mu_B|$.
The dimension is $3$ (structural for all $q \ge 3$).

\subsection{$\mu_C$ piece is zero by $j$-formula leakage}

\begin{lemma}[$\mu_C$ piece vanishes]\label{lem:muC-zero}
For $q \ge 1$, $P^{(\mu_C)}_{j_0-1, n-2}\, V_{\mu_C} = 0$.
\end{lemma}

\begin{proof}
Since $\ell(\mu_C) = q + 2 = \ell(\lambda)$, the sharp index for
$\mu_C$ is $j_0(\mu_C) = q + 3 = j_0(\lambda)$. Our inner iteration
ranges $j_0 - 1$ to $n - 2 = q + 5$, one step beyond the sharp index
of $\mu_C$:
\[
   P^{(\mu_C)}_{j_0-1, n-2} V_{\mu_C}
   \;=\; R'^{(\mu_C)}_{j_0-1}\, P^{(\mu_C)}_{j_0, n-2}\, V_{\mu_C}
   \;=\; R'^{(\mu_C)}_{q+2}\,\big(B^{(\mu_C)}_{j_0(\mu_C)} V_{\mu_C}\big).
\]
By Theorem~\ref{thm:muC}, $B^{(\mu_C)}_{j_0(\mu_C)} V_{\mu_C}
\subseteq E_1^-|_{V_{\mu_C}}$. The operator $R'^{(\mu_C)}_{q+2}
= (T_{q+1}{+}1)(T_q{+}1)\cdots(T_2{+}1)(T_1{+}1)$ has rightmost factor
$(T_1{+}1)$, which is applied first to any input vector. On
$E_1^-$, $T_1$ acts as $-1$, so $(T_1{+}1)|_{E_1^-} = 0$. Hence
$R'^{(\mu_C)}_{q+2}\, E_1^-|_{V_{\mu_C}} = 0$, giving the claim.\qed
\end{proof}

\subsection{$S$ piece vanishes by the same-row-dies framework}

\begin{lemma}[$S$ piece vanishes]\label{lem:S-zero}
For $q \ge 1$, $P_{j_0, n-1}\, S = 0$.
\end{lemma}

\begin{proof}
We apply Theorem~\ref{thm:srd-framework} (general same-row-dies
framework) with $\widetilde\lambda = \lambda$ and $i = 2$:
\begin{itemize}[leftmargin=2em,topsep=0.2em,itemsep=0.1em]
\item Same-row pair at row $2$: $\lambda_2 = 3 \ge 2$ and
$\lambda_3 = 1 \le \lambda_2 - 2 = 1$.~\checkmark
\item $\lambda_2 = 3 \ge 3$, so $\ell(\nu_2) = \ell(\lambda)$ (since
removing two cells $(2, 2), (2, 3)$ leaves $(2, 1)$ as the row-$2$
cell, preserving row $2$).~\checkmark
\item $\nu_2 = (4, 1^{q+1})$ is a hook on $n - 2 = q + 5 \ge 8$
letters. By Theorem~\ref{thm:hook-j} (Shift Lemma all-hooks)
with $k = 4$ and $m = q + 1$,
$B^{(\nu_2)}_{j_0(\nu_2)} V_{\nu_2} \subseteq E_1^-|_{V_{\nu_2}}$.~\checkmark
\end{itemize}
Hence $P_{j_0, n-1}\, S = R'_{j_0} R'_{j_0 + 1} \cdots R'_{n-1}\, S
= 0$.\qed
\end{proof}

\section{Main theorem: $j$-formula and upper bound}\label{sec:main}

\begin{theorem}[$j$-formula at sharp index]\label{thm:main-j}
For $\lambda = (4, 3, 1^q)$ at every $q \ge 3$,
\[
   B^{(\lambda)}_{j_0} V_\lambda \;\subseteq\; E_1^-\!\mid_{V_\lambda}.
\]
\end{theorem}

\begin{proof}
By the four-branch reduction
(Theorem~\ref{thm:reduction}) and Lemmas~\ref{lem:muC-zero}
and~\ref{lem:S-zero},
\[
   B^{(\lambda)}_{j_0} V_\lambda
   \;=\;
   (T_{j_0-1}{+}1)(T_{j_0}{+}1)\cdots(T_{n-2}{+}1)
   \Big[ \Phi_A\big(B^{(\mu_A)}_{j_0(\mu_A)} V_{\mu_A}\big)
        + \Phi_B\big(B^{(\mu_B)}_{j_0(\mu_B)} V_{\mu_B}\big) \Big].
\]
By Theorems~\ref{thm:muA}(1) and~\ref{thm:muB}(2) for $q \ge 3$, both
inner images lie in $E_1^-$ of the respective $V_{\mu_X}$. By
Theorem~\ref{thm:Phi}(4), $\Phi_X$ identifies $E_1^-|_{V_{\mu_X}}$
with $\Phi_X(V_{\mu_X}) \cap E_1^-|_{V_\lambda}$. Hence
\[
   \Phi_A\big(B^{(\mu_A)}_{j_0(\mu_A)} V_{\mu_A}\big)
   \;\subseteq\; E_1^-\!\mid_{V_\lambda},
   \qquad
   \Phi_B\big(B^{(\mu_B)}_{j_0(\mu_B)} V_{\mu_B}\big)
   \;\subseteq\; E_1^-\!\mid_{V_\lambda}.
\]
Finally, each outer factor $(T_j{+}1)$ for $j \in \{j_0 - 1, j_0,
\ldots, n - 2\}$ has index $j \ge j_0 - 1 = q + 2 \ge 5$ (since
$q \ge 3$). Each such $T_j$ commutes with $T_1$ (index gap
$j - 1 \ge 4$), so each $(T_j{+}1)$ preserves $E_1^-|_{V_\lambda}$.
Multiplying these together still preserves $E_1^-$.
\qed
\end{proof}

\begin{corollary}[Rank-zero]\label{cor:rank-zero}
$\Pi^{S_n}|_{V_\lambda} = B^{(\lambda)}_2 V_\lambda = 0$.
\end{corollary}

\begin{proof}
From Theorem~\ref{thm:main-j}, $B^{(\lambda)}_{j_0} V_\lambda \subseteq
E_1^-$. Iterating downward: $B^{(\lambda)}_{j_0 - 1} V_\lambda
= R'_{j_0-1}\, B^{(\lambda)}_{j_0} V_\lambda \subseteq R'_{j_0-1}\, E_1^-
= 0$, since $R'_{j_0-1}$ ends with $(T_1{+}1)$ on the right (and
$j_0 - 1 \ge 2$). Successive iterations preserve $0$, giving
$B^{(\lambda)}_2 V_\lambda = 0$.\qed
\end{proof}

\begin{theorem}[Upper bound]\label{thm:main-upper}
For $\lambda = (4, 3, 1^q)$ at every $q \ge 3$,
\[
   \dim B^{(\lambda)}_{j_0} V_\lambda \;\le\; 5.
\]
The upper bound is structural for $q \ge 5$ and mixed
structural-computational for $q \in \{3, 4\}$.
\end{theorem}

\begin{proof}
From the proof of Theorem~\ref{thm:main-j},
\[
   B^{(\lambda)}_{j_0} V_\lambda
   \;=\;
   (T_{j_0-1}{+}1)\cdots(T_{n-2}{+}1)
   \Big[ \Phi_A\big(B^{(\mu_A)}_{j_0(\mu_A)} V_{\mu_A}\big)
        + \Phi_B\big(B^{(\mu_B)}_{j_0(\mu_B)} V_{\mu_B}\big) \Big].
\]
Each $\Phi_X$ is injective, so
$\dim \Phi_X(B^{(\mu_X)}_{j_0(\mu_X)} V_{\mu_X}) = \dim B^{(\mu_X)}_{j_0(\mu_X)} V_{\mu_X}$.
The sum of dimensions in the bracket is therefore at most
\[
   \dim B^{(\mu_A)}_{j_0(\mu_A)} V_{\mu_A}
   + \dim B^{(\mu_B)}_{j_0(\mu_B)} V_{\mu_B}
   \;=\; 2 + 3 \;=\; 5,
\]
using Theorems~\ref{thm:muA}(2) and~\ref{thm:muB}(1).
Multiplying by the outer factors $(T_{j}{+}1)$ cannot increase the
dimension. Hence $\dim B^{(\lambda)}_{j_0} V_\lambda \le 5$.

For $q \ge 5$, both summand dimensions are structural; for
$q \in \{3, 4\}$, $\dim B^{(\mu_A)} = 2$ at $\mu_A = (3, 3, 1^{q-1})$
with $q - 1 \in \{2, 3\}$ is computational (SRD-331,
Remark~\ref{thm:muA}).\qed
\end{proof}

\section{Computational verification}\label{sec:comp}

\paragraph{Direct dim test for $\lambda = (4, 3, 1^q)$.} Using
\texttt{compute\_dim\_B.py} (May-13 multi-corner paper's script),
mod $P = 100\,003$ at $q_{\text{spec}} = 11$:

\begin{center}
\begin{tabular}{lcc}
\toprule
$\lambda$ & $\dim V_\lambda$ & $\dim B^{(\lambda)}_{j_0} V_\lambda$ \\
\midrule
$(4, 3, 1, 1, 1)$ \quad ($q = 3$, $n = 10$) & $525$ & $5$ \\
$(4, 3, 1, 1, 1, 1)$ \quad ($q = 4$, $n = 11$) & $1100$ & $5$ \\
\bottomrule
\end{tabular}
\end{center}

Both match the conjectural $f^{(3, 2)} = 5$, agreeing with the
structural upper bound and the May-13 multi-corner closed form.

\paragraph{Same-row part dimension.} For the $S$-piece:
$\dim S = f^{\nu_2} = f^{(4, 1^{q+1})} = \binom{q+4}{3}$.
At $q = 3$: $\dim S = 35$. At $q = 4$: $\dim S = 56$. Despite
this large dimension, Lemma~\ref{lem:S-zero} forces
$P_{j_0, n-1} S = 0$.

\section{Open problems and follow-up}\label{sec:open}

\paragraph{Structural lower bound: $\dim B \ge 5$.}
The upper bound $\dim B \le 5$ is matched by the computational
lower bound $\dim B \ge 5$ at $q \in \{3, 4\}$. A structural lower
bound would construct $5$ explicit vectors $F_{X, Y}$ in
$B^{(\lambda)}_{j_0} V_\lambda$ (indexed by an SYT of $\mu_A$ or
$\mu_B$ combined with a corner-pair label) and separate them via
position-of-$n$ / position-of-$(n-1)$ projections, analogous to the
May-12-Eve3 lower-bound argument for $(4, 2, 1^*)$. We defer this to
a follow-up paper.

\paragraph{Dim formula at $q \in \{3, 4\}$, fully structural.}
Theorem~\ref{thm:main-upper}'s upper bound at $q \in \{3, 4\}$ uses
the computational dim $= 2$ for $\mu_A = (3, 3, 1^{q-1})$ at
$q - 1 \in \{2, 3\}$. A structural lower bound on $\mu_A$
(currently open below $q - 1 = 4$) would propagate the structural
upper bound on $\lambda$ to $q = 3, 4$ as well.

\paragraph{Beyond $(4, 3, 1^*)$: the $(k, 3, 1^*)$ family.}
The proof of Lemma~\ref{lem:S-zero} (same-row-dies for $\lambda$)
applies to any $(k, 3, 1^q)$ at $k \ge 4$ with $\lambda_2 = 3$
and $\lambda_3 = 1$: the sub-shape $\nu_2 = (k, 1^{q+1})$ is a
hook and Theorem~\ref{thm:hook-j} applies. The four-branch
reduction (Theorem~\ref{thm:reduction}) and the $\mu_X$ identifications
analogously give the upper bound
$\dim B^{(\lambda)}_{j_0} V_\lambda \le f^{(k-1, 2)}$
\emph{provided} the recursive sub-shape inputs
($\mu_A = (k-1, 3, 1^{q-1})$, $\mu_B = (k, 2, 1^{q-1})$,
$\mu_C = (k-1, 2, 1^q)$) themselves have known dim/$j$-formulas.
For $k = 5$: $\mu_A = (4, 3, 1^{q-1})$ is the family of this paper
(at one-smaller $q$), giving a recursive bootstrap upward.

\paragraph{The general program.}
Combining the same-row-dies framework with the multi-corner recursion
gives a roadmap for proving the closed form
$\dim B^{(\lambda)}_{j_0} V_\lambda = f^{\lambda^{(\ge 2)}}$
structurally across the stable regime: each step adds a new shape
class to the proven $j$-formulas, which in turn unlocks new
sub-shape inputs via the same-row-dies framework. This paper closes
$(4, 3, 1^*)$; the next natural targets are $(5, 3, 1^*)$ and
$(5, 4, 1^*)$, which appear within reach using the same toolkit.

\section{Honest status}\label{sec:honest}

\paragraph{What we proved.}
For $\lambda = (4, 3, 1^q)$ at every $q \ge 3$:
\begin{enumerate}[leftmargin=2em,topsep=0.2em,itemsep=0.1em]
\item The $j$-formula at sharp index:
$B^{(\lambda)}_{j_0} V_\lambda \subseteq E_1^-|_{V_\lambda}$, hence
rank-zero $\Pi^{S_n}|_{V_\lambda} = 0$
(Theorem~\ref{thm:main-j}, Corollary~\ref{cor:rank-zero}).
\item The upper bound $\dim B^{(\lambda)}_{j_0} V_\lambda \le 5$
(Theorem~\ref{thm:main-upper}), structurally for $q \ge 5$ and
mixed structural-computational for $q \in \{3, 4\}$.
\end{enumerate}

\paragraph{What we did not prove.}
The matching structural lower bound $\dim \ge 5$.
Computationally, $\dim = 5$ at $q \in \{3, 4\}$ from
\texttt{compute\_dim\_B.py}.

\paragraph{What this closes.}
This is the first $r = 2$ same-row family beyond $(4, 2, 1^*)$ where
the structural $j$-formula at sharp index is established. It
demonstrates that the May-13 same-row-dies framework, combined with
the prior $j$-formulas on hooks and small non-hook shapes, propagates
across an infinite-$q$ family. The proof method is robust and points
toward $(5, 3, 1^*)$, $(k, 3, 1^*)$, and recursive bootstrapping up
the multi-corner ladder.

\end{document}
